\subsection{Witt polynomials}

For every integer \(n \geqslant 0\), we call the element \(\Phi_n\) of \(\mathbf{Z}[X_0, ..., X_n]\) the \(n\)-th Witt polynomial, defined by

\[
\Phi_ {n} \left(\mathrm{X} _ {0}, \dots , \mathrm{X} _ {n}\right) = \sum_ {i = 0} ^ {n} p ^ {i} \mathrm{X} _ {i} ^ {p ^ {n - i}} = \mathrm{X} _ {0} ^ {p ^ {n}} + p \mathrm{X} _ {1} ^ {p ^ {n - 1}} + \dots + p ^ {n} \mathrm{X} _ {n}.\tag{1}
\]

One evidently has \(\Phi_{0}=X_{0}\) and the recurrence relations

\begin{enumerate}
\def\labelenumi{(\arabic{enumi})}
\setcounter{enumi}{1}
\tightlist
\item
\end{enumerate}

\[
\Phi_ {n + 1} \left(\mathrm{X} _ {0}, \dots , \mathrm{X} _ {n + 1}\right) = \Phi_ {n} \left(\mathrm{X} _ {0} ^ {p}, \dots , \mathrm{X} _ {n} ^ {p}\right) + p ^ {n + 1} \mathrm{X} _ {n + 1}\tag{3}
\]

\[
\Phi_ {n + 1} \left(\mathrm{X} _ {0}, \dots , \mathrm{X} _ {n + 1}\right) = \mathrm{X} _ {0} ^ {p ^ {n + 1}} + p \Phi_ {n} \left(\mathrm{X} _ {1}, \dots , \mathrm{X} _ {n + 1}\right).
\]

When one assigns \(X_{i}\) the weight \(p^{i}\), the polynomial \(\Phi_{n}\) is isobaric of weight \(p^{n}\) (A, IV, p.3).

PROPOSITION 1. --- Let A be a filtered ring and \((\mathbf{J}_n)_{n\in \mathbb{Z}}\) its filtration. We assume that \(\mathbf{J}_0 = \mathbf{A}\) and \(p\cdot 1_{\mathbf{A}}\in \mathbf{J}_{1}\). Let m and n be integers such that \(m\geqslant 1\) and \(n\geqslant 0\), and let \(a_0,\dots,a_n,b_0,\dots,b_n\) be elements of A.

\begin{enumerate}
\def\labelenumi{\alph{enumi})}
\tightlist
\item
  If \(a_i \equiv b_i \mod J_m\) for \(0 \leqslant i \leqslant n\), then we have
\end{enumerate}

\[
\Phi_ {i} (a _ {0}, \dots , a _ {i}) \equiv \Phi_ {i} (b _ {0}, \dots , b _ {i}) \bmod \mathrm{J} _ {m + i} \quad \text{for} \quad 0 \leqslant i \leqslant n.
\]

\begin{enumerate}
\def\labelenumi{\alph{enumi})}
\setcounter{enumi}{1}
\tightlist
\item
  Suppose that, for every integer \(k \geqslant 1\) and for all \(x \in \mathbf{A}\), the relation \(p x \in \mathbf{J}_{k+1}\) implies \(x \in \mathbf{J}_k\). If one has \(\Phi_i(a_0, ..., a_i) \equiv \Phi_i(b_0, ..., b_i) \mod \mathbf{J}_{m+i}\) for \(0 \leqslant i \leqslant n\), then we have \(a_i \equiv b_i \mod \mathbf{J}_m\) for \(0 \leqslant i \leqslant n\).
\end{enumerate}

Lemma 1. --- If x and y are two elements of A congruent modulo \(J_{m}\), one has

\[
x ^ {p ^ {n}} \equiv y ^ {p ^ {n}} \bmod J _ {m + n}.
\]

By induction on \(n\) , we reduce to the case where \(n=1\) . Let P denote the polynomial \(\sum_{i=0}^{p-1} X^i Y^{p-1-i}\) of Z{[}X, Y{]}. In view of the hypothesis made on \(x\) and \(y\) , one has P(x, y) \(\equiv\) P(x, x) \(\equiv\) p, \(x^{p-1}\) mod. J \(_m\) . But J \(_m\) + p.A \(\subset\) J \(_1\) , hence P(x, y) \(\in\) J \(_1\) . Finally, \(x^p - y^p = (x - y)\) P(x, y) belongs to J \(_m\) J \(_1\) \(\subset\) J \(_{m+1}\) .

Let us prove a) by induction on n. The case n = 0 is immediate. Suppose \(n \geqslant 1\). Under the hypotheses of a), we have

\begin{enumerate}
\def\labelenumi{(\arabic{enumi})}
\setcounter{enumi}{3}
\tightlist
\item
\end{enumerate}

\(a_{i}^{p}\equiv b_{i}^{p}\bmod \mathrm{J}_{m + 1}\quad \text{for}\quad 0\leqslant i\leqslant n - 1\quad \text{by Lemma } 1.\)

\[
\Phi_ {n - 1} (a _ {0} ^ {p}, \dots , a _ {n - 1} ^ {p}) \equiv \Phi_ {n - 1} (b _ {0} ^ {p}, \dots , b _ {n - 1} ^ {p}) \pmod{J _ {m + n}}\tag{5}
\]

by the induction hypothesis applied to the elements \(a_0^p, \ldots, a_{n-1}^p, b_0^p, \ldots, b_{n-1}^p\) of A, and

\[
\Phi_ {n} (a _ {0},..., a _ {n}) - p ^ {n}. a _ {n} \equiv \Phi_ {n} (b _ {0},..., b _ {n}) - p ^ {n}. b _ {n} \pmod{J _ {m + n}}\tag{6}
\]

by formulas (2) and (5). Since \(a_{n}-b_{n}\) belongs to \(\mathbf{J}_{m}\), the element \(p^{n}a_{n}-p^{n}b_{n}\) belongs to \(\mathbf{J}_{m+n}\), and from (6) we deduce the congruence

\[
\Phi_ {n} (a _ {0}, \dots , a _ {n}) \equiv \Phi_ {n} (b _ {0}, \dots , b _ {n}) \pmod{J _ {m + n}},
\]

whence a).

We prove \(b)\) by induction on \(n\). The case \(n=0\) is immediate. Suppose \(n\geqslant1\). Under the hypotheses of \(b)\), we have \(a_{i}\equiv b_{i}\) mod. \(\mathbf{J}_{m}\) for \(0\leqslant i\leqslant n-1\) by the induction hypothesis, and we deduce as before the congruences (4), (5), and (6). But by hypothesis \(\Phi_{n}(a_{0},...,a_{n})\) and \(\Phi_{n}(b_{0},...,b_{n})\) are congruent mod. \(\mathbf{J}_{m+n}\), and we therefore have \(p^{n}(a_{n}-b_{n})\in\mathbf{J}_{m+n}\). Since the relation \(p x\in\mathbf{J}_{k+1}\) implies \(x\in\mathbf{J}_{k}\) for every \(x\in\mathbf{A}\) and every \(k\geqslant1\), one has \(a_{n}-b_{n}\in\mathbf{J}_{m}\), which completes the proof.

